\documentclass[11pt]{article}
\usepackage[utf8]{inputenc}
\usepackage[margin=1in]{geometry}
\usepackage{amsmath}
\usepackage{graphicx}
\usepackage{booktabs}
\usepackage{hyperref}
\usepackage{natbib}
\usepackage{lineno}
\usepackage{xcolor}
\usepackage{multirow}
\usepackage{float}

\linenumbers

\title{Auditory Cortex Conveys Non-Topographic Sound Location Information to Primary Visual Cortex}
\author{
  Jennifer L.\ Bizley$^{1}$ \and
  Andrew J.\ King$^{2}$
}
\date{
  $^{1}$Ear Institute, University College London, London, UK \\
  $^{2}$Department of Physiology, Anatomy and Genetics, University of Oxford, Oxford, UK
}

\begin{document}

\maketitle

\begin{abstract}
Auditory cortex (AC) sends direct projections to primary visual cortex (V1) that are thought to mediate cross-modal interactions, but whether these projections carry spatial auditory information---and whether this information is topographically organized to match V1's retinotopic map---was unknown. Using dual-color two-photon calcium imaging in head-fixed mice, we simultaneously recorded AC axon boutons (GCaMP6s) in V1 layer 1 and postsynaptic V1 neurons (jRGECO1a) in layer 2/3, while presenting spatially controlled auditory and visual stimuli from a custom 39-position speaker-LED array. AC boutons in V1 carried significant information about sound source location: population decoding of azimuth achieved $72.4 \pm 3.1$\% accuracy (chance: $7.7$\%; $p < 0.001$, $n = 7$ mice) and elevation achieved $58.3 \pm 4.2$\% accuracy (chance: $33.3$\%; $p < 0.001$). However, auditory spatial receptive fields of AC boutons were not topographically aligned with the visual receptive fields of their postsynaptic V1 neurons (correlation $r = 0.04$, $p = 0.71$). AC boutons spanned a wider spatial range than V1's visual field, with substantial fractions tuned to ipsilateral sound locations beyond V1's visual coverage. In contrast, secondary visual area lateral (V2L) inputs to V1 showed strong visual retinotopy but minimal auditory spatial coding (azimuth decoding: $12.8 \pm 2.9$\%, $p = 0.18$ vs.\ shuffle). Audiovisual congruency experiments revealed that sound-evoked modulation of V1 responses did not depend on spatial alignment between sound and light ($F(4,28) = 0.42$, $p = 0.79$). These results demonstrate that AC provides V1 with a non-topographic representation of auditory space that may serve functions beyond spatially congruent audiovisual binding.
\end{abstract}

\section{Introduction}

Binding spatiotemporally congruent multisensory stimuli is fundamental to coherent perception. In the superior colliculus (SC), auditory and visual spatial maps are in topographic register, providing a neural substrate for matching stimuli across modalities based on spatial location \citep{wallace1996, meredith1987, stein2008}. Whether a similar topographic alignment exists in cortex---particularly between auditory cortex projections and the retinotopic map in primary visual cortex---has been an open question \citep{ghazanfar2006, scholl2021}.

AC sends direct projections to V1 across species \citep{bizley2007, iurilli2012, ibrahim2016}, and these projections are concentrated in layer 1, where they contact the apical dendrites of pyramidal neurons in layers 2/3 \citep{ibrahim2016}. Neurons in AC exhibit spatial tuning to sound source location through population coding rather than topographic maps \citep{middlebrooks1994, stecker2005}. However, it was unknown whether the AC-to-V1 projection specifically relays this spatial information, and if so, whether the auditory spatial representation is organized to match V1's retinotopic map.

Here, using a custom 39-position speaker-LED array and dual-color two-photon imaging, we simultaneously characterize the auditory spatial properties of AC inputs and the visual receptive fields of their postsynaptic V1 targets. We complement this with V2L-to-V1 input characterization as a control pathway and test whether audiovisual modulation in V1 depends on spatial congruency.

\section{Results}

\subsection{Experimental design and imaging}

We injected AAV-GCaMP6s into AC ($n = 7$ mice) or V2L ($n = 3$ mice) and AAV-jRGECO1a into V1, enabling simultaneous imaging of presynaptic boutons in V1 layer 1 (green channel) and postsynaptic somata in V1 layer 2/3 (red channel). A 39-position array (13 azimuth columns $\times$ 3 elevation rows, 10\textdegree\ azimuth steps, 20\textdegree\ elevation steps) presented white-noise bursts (2--20 kHz, 1 s) and flashing LEDs (Table~\ref{tab:setup}).

\begin{table}[H]
\centering
\caption{Experimental parameters for the stimulus array and imaging setup.}
\label{tab:setup}
\begin{tabular}{lcc}
\toprule
\textbf{Parameter} & \textbf{Value} & \textbf{Units} \\
\midrule
Array positions & 39 & --- \\
Azimuth columns & 13 & --- \\
Elevation rows & 3 & --- \\
Azimuth step & 10 & degrees \\
Elevation step & 20 & degrees \\
Ipsilateral coverage & 20 & degrees \\
Contralateral coverage & 100 & degrees \\
Auditory stimulus & 2--20 kHz WN & --- \\
Stimulus duration & 1.0 & seconds \\
AC-injected mice & 7 & --- \\
V2L-injected mice & 3 & --- \\
Imaging sessions (AC) & 42 & --- \\
Imaging sessions (V2L) & 15 & --- \\
\bottomrule
\end{tabular}
\end{table}

\subsection{Modality selectivity of AC vs.\ V2L inputs}

AC boutons in V1 were predominantly auditory-responsive, while V2L boutons were predominantly visual-responsive (Table~\ref{tab:modality}).

\begin{table}[H]
\centering
\caption{Fraction of responsive boutons by modality in AC and V2L inputs (mean $\pm$ SD across mice). Two-sided paired $t$-tests comparing auditory vs.\ visual fractions within each input type.}
\label{tab:modality}
\begin{tabular}{lcccc}
\toprule
\textbf{Input} & \textbf{Auditory (\%)} & \textbf{Visual (\%)} & \textbf{Both (\%)} & \textbf{$p$ (aud.\ vs.\ vis.)} \\
\midrule
AC boutons & $68.4 \pm 8.2$ & $12.3 \pm 4.1$ & $8.7 \pm 3.2$ & $< 0.001$ \\
V2L boutons & $14.2 \pm 5.3$ & $71.8 \pm 9.1$ & $6.4 \pm 2.8$ & $< 0.001$ \\
\bottomrule
\end{tabular}
\end{table}

\subsection{Population decoding of sound location from AC inputs}

Maximum-likelihood decoding of azimuth from AC bouton populations achieved $72.4 \pm 3.1$\% accuracy, far exceeding chance (7.7\%) and shuffled controls ($8.1 \pm 1.4$\%). Elevation decoding reached $58.3 \pm 4.2$\% (chance: 33.3\%). V2L inputs showed near-chance auditory decoding (Table~\ref{tab:decoding}).

\begin{table}[H]
\centering
\caption{Sound location decoding accuracy ($\uparrow$ higher is better) from AC and V2L bouton populations in V1 (mean $\pm$ SD). Two-sided paired $t$-tests comparing real vs.\ shuffled data.}
\label{tab:decoding}
\begin{tabular}{lcccccc}
\toprule
\textbf{Dimension} & \textbf{AC (\%)} & \textbf{AC shuffle (\%)} & \textbf{$p$ (AC)} & \textbf{V2L (\%)} & \textbf{V2L shuffle (\%)} & \textbf{$p$ (V2L)} \\
\midrule
Azimuth & $72.4 \pm 3.1$ & $8.1 \pm 1.4$ & $< 0.001$ & $12.8 \pm 2.9$ & $8.3 \pm 1.2$ & $0.182$ \\
Elevation & $58.3 \pm 4.2$ & $34.1 \pm 2.1$ & $< 0.001$ & $36.2 \pm 3.8$ & $33.7 \pm 1.9$ & $0.341$ \\
\bottomrule
\end{tabular}
\end{table}

\begin{table}[H]
\centering
\caption{Mean decoding error (degrees, $\downarrow$ lower is better) between decoded and actual speaker position for AC and V2L inputs.}
\label{tab:decoding_error}
\begin{tabular}{lccc}
\toprule
\textbf{Dimension} & \textbf{AC error (deg)} & \textbf{V2L error (deg)} & \textbf{$p$ (AC vs.\ V2L)} \\
\midrule
Azimuth & $14.2 \pm 2.8$ & $41.3 \pm 6.1$ & $< 0.001$ \\
Elevation & $9.8 \pm 2.1$ & $18.4 \pm 3.7$ & $0.003$ \\
\bottomrule
\end{tabular}
\end{table}

\subsection{AC inputs are not retinotopically organized}

The peak azimuth of auditory spatial tuning in AC boutons showed no correlation with the visual RF center of their postsynaptic V1 neurons (Table~\ref{tab:retinotopy}).

\begin{table}[H]
\centering
\caption{Correlation between auditory RF peak (AC boutons) and visual RF center (V1 somata) across imaging fields (Pearson correlation with 95\% CI). No significant retinotopic alignment was detected.}
\label{tab:retinotopy}
\begin{tabular}{lccccc}
\toprule
\textbf{Comparison} & \textbf{$r$} & \textbf{95\% CI} & \textbf{$p$} & \textbf{$n$ (fields)} \\
\midrule
Aud.\ azimuth vs.\ Vis.\ azimuth & $0.04$ & $[-0.18, 0.26]$ & $0.712$ & 42 \\
Aud.\ elevation vs.\ Vis.\ elevation & $0.08$ & $[-0.14, 0.29]$ & $0.481$ & 42 \\
V2L vis.\ azimuth vs.\ V1 vis.\ azimuth & $0.87$ & $[0.71, 0.94]$ & $< 0.001$ & 15 \\
\bottomrule
\end{tabular}
\end{table}

\subsection{Audiovisual congruency does not modulate V1 responses}

V1 neuron responses to a fixed LED in the visual RF were modulated by concurrent sound, but this modulation did not depend on spatial alignment (Table~\ref{tab:congruency}).

\begin{table}[H]
\centering
\caption{Audiovisual congruency experiment. Mean V1 response amplitude ($\Delta F/F$, $\times 10^{-2}$) to LED paired with sound at varying spatial offsets. One-way repeated-measures ANOVA across speaker positions.}
\label{tab:congruency}
\begin{tabular}{lcccccc}
\toprule
\textbf{Speaker offset (deg)} & $-40$ & $-20$ & $0$ & $+20$ & $+40$ & \textbf{ANOVA} \\
\midrule
$\Delta F/F$ ($\times 10^{-2}$) & $3.21 \pm 0.42$ & $3.18 \pm 0.39$ & $3.24 \pm 0.44$ & $3.19 \pm 0.41$ & $3.15 \pm 0.38$ & $F(4,28)=0.42$ \\
& & & & & & $p = 0.794$ \\
\bottomrule
\end{tabular}
\end{table}

\section{Discussion}

Our results reveal that auditory cortex provides primary visual cortex with a rich representation of sound source location, but this representation is not organized topographically to match V1's retinotopic map. Three key findings emerge from our experiments.

First, AC axon boutons in V1 layer 1 carry substantial auditory spatial information, with population decoding achieving $72.4$\% accuracy for azimuth---far above chance---confirming that the AC-to-V1 projection is a significant channel for spatial auditory information. This extends prior work showing that V1-projecting AC neurons encode specific sound features \citep{bizley2007, king2009}.

Second, the absence of retinotopic alignment between auditory spatial RFs and visual RFs ($r = 0.04$) stands in striking contrast to the superior colliculus, where auditory and visual maps are in spatial register \citep{wallace1996, stein2008, knudsen2002}. This suggests that if cortical audiovisual binding depends on spatial congruency, it must use a mechanism other than topographic map alignment.

Third, the lack of spatial congruency dependence in audiovisual modulation of V1 responses ($p = 0.79$) suggests that AC inputs to V1 may serve functions beyond spatially congruent stimulus binding, such as providing a global alerting signal that enhances visual processing regardless of sound location \citep{iurilli2012, ibrahim2016}.

The broad spatial coverage of AC inputs---spanning well into ipsilateral space beyond V1's visual field---further argues against a role in fine-grained spatial matching and instead suggests that AC conveys a panoramic auditory spatial representation that could serve orienting or attentional functions \citep{middlebrooks1994, ghazanfar2006}.

\section{Methods}

\subsection{Animals and viral injections}

Head-fixed C57BL/6J mice received injections of AAV-Syn-GCaMP6s \citep{chen2012} into AC ($n = 7$) or V2L ($n = 3$) and AAV-Syn-jRGECO1a \citep{dana2016} into V1. Chronic cranial windows were implanted over V1. All procedures were approved by the institutional animal care committee.

\subsection{Stimulus delivery}

A custom array of 39 speaker-LED pairs (13 columns $\times$ 3 rows) delivered white-noise bursts (2--20 kHz, 70 dB SPL, 1 s duration) and flashing white LEDs. Extended azimuth experiments used a single speaker moved from $-90$\textdegree\ to $+90$\textdegree\ in 20\textdegree\ steps with broadband noise (2--80 kHz).

\subsection{Two-photon calcium imaging}

Dual-color volumetric two-photon imaging was performed at 920 nm (GCaMP6s) and 1040 nm (jRGECO1a). AC/V2L boutons were imaged in V1 layer 1; V1 somata in layer 2/3.

\subsection{Analysis}

Spatial modulation index (SMI) quantified selectivity. Maximum-likelihood decoding used leave-one-out cross-validation. Shuffled controls randomized trial labels. Retinotopic alignment was assessed by Pearson correlation between bouton auditory RF peaks and soma visual RF centers.

\section{Conclusion}

Auditory cortex projections to V1 convey accurate sound location information through a non-topographic representation that spans both ipsilateral and contralateral auditory space. The absence of retinotopic organization and spatial congruency effects on audiovisual modulation suggests that these inputs serve functions beyond spatially matched audiovisual binding, potentially including cross-modal attentional or contextual modulation of visual processing.

\bibliographystyle{plainnat}
\bibliography{references}

\end{document}
